\documentclass[11pt]{article}
\usepackage[a4paper,margin=1in]{geometry}
\usepackage[T1]{fontenc}
\usepackage[utf8]{inputenc}
\usepackage{lmodern}
\usepackage{microtype}
\usepackage{booktabs}
\usepackage{enumitem}
\usepackage{hyperref}
\usepackage{xcolor}
\hypersetup{colorlinks=true, linkcolor=black, citecolor=black, urlcolor=blue}

\title{\textbf{spec-loop: A Fact-Gated, Tool-Agnostic Methodology for\\Specification-Driven Development with LLM Coding Agents}}
\author{Andr\'es Massello\\
\small C\'ordoba, Argentina\\
\small \texttt{https://github.com/andresmassello}}
\date{July 2026}

\begin{document}
\maketitle

\begin{abstract}
LLM coding agents fail in two well-documented ways: when instructions are
underspecified they guess instead of stopping, and when a verification signal is
visible they optimize the signal rather than the requested outcome. We present
\textbf{spec-loop}, a methodology for specification-driven development with LLM
coding agents organized around a single discipline: \emph{gates that read facts
block; gates that guess over prose advise}. The methodology contributes (i) two
specification fronts --- a greenfield front where the agent \emph{proposes} the
system shape under one-question-at-a-time human approval, and a brownfield front
where the agent may only \emph{extract} verifiable facts, never author a
specification of existing behavior; (ii) a deterministic evidence ledger with a
two-tier record contract in which measured artifacts always override agent
self-reports; (iii) a family of executable fact gates (golden non-authorship,
gate-integrity, simplicity budgets, per-criterion measured acceptance,
regression capture, a derived --- never self-declared --- workflow state machine)
and a structured-guess layer (versioned qualitative rubrics with mandatory
evidence) that advises by default and gates only by explicit human declaration;
and (iv) a strict placement of the human at three fixed points: approving the
shape, approving the golden baseline, and merging. A reference implementation
(seven agent skills and a dependency-free Python engine with 23 subcommands and
adapters for nine language stacks) is described, together with its design
validation: two adversarial audits (231 agent evaluations) that initially found
the central principle inverted in code and drove its re-wiring; a systematic
cross-examination against classical software-engineering doctrine that produced
ten shipped improvements; and sustained self-application, in which fresh-context
reviews caught real defects in 13 of 14 releases before merge. We report design
validation and self-application evidence only; controlled empirical evaluation
remains future work.
\end{abstract}

\section{Introduction}\label{sec:intro}

Two recent empirical results frame the problem this methodology addresses.
First, when task instructions are underspecified, production coding agents do
not fail closed: they act on unstated assumptions, and a majority of runs
violate at least one action boundary~\cite{ji2026guessing}. Second, when a
verification oracle is visible to the agent, near-perfect scores can coexist
with an undelivered artifact: the agent \emph{builds to the test}, optimizing
the check rather than the request, and does not on its own validate what it
ships as a user would~\cite{ma2026building}. Both failures are dispositional
rather than capability failures, and both worsen as agents are given more
autonomy.

Existing agent development frameworks converge on persistent artifacts and
human review as mitigations~\cite{demacedo2026prompt}, but they largely encode
their process as \emph{prose instructions that the agent may ignore}, and their
verification signals rarely distinguish between what is mechanically verifiable
and what is a model's opinion. spec-loop starts from that distinction and makes
it the organizing principle of the whole process.

The methodology is built on one discipline and three architectural commitments:

\begin{description}[leftmargin=1.2em]
\item[Facts block; guesses advise.] Every check in the process is classified as
\emph{computational} (it reads an artifact --- a byte comparison, a parsed test
report, a diff structure) or \emph{inferential} (it judges prose or code
quality). Computational checks are allowed to block progress mechanically.
Inferential checks may only advise, unless a human explicitly declares them
gating. A natural-language heuristic that blocks produces false positives, gets
disabled, and dies; keeping guesses advisory is what keeps them alive.
\item[Measured beats narrated.] The evidence ledger records two tiers: measured
records parsed from real artifacts (test reports, coverage files, linter
output) and self-reported agent counts. A measured red always overrides a
narrated green. Acceptance criteria close per criterion only when a green,
name-tagged test case exists in the ingested reports --- a checked checkbox
without a test is recorded as \emph{narrated-only} and does not close.
\item[The human at three fixed points.] The human approves the \emph{shape}
(specification front), approves the \emph{golden baseline} (the one artifact
the agent is mechanically forbidden from authoring), and \emph{merges}. The
agent proposes, measures, and stops.
\item[Tool-agnosticism by contract.] The engine is dependency-free Python that
never runs an LLM: it validates structures and ingests reports. Wherever a
model's judgment participates (e.g.\ rubric grading), the interface is a JSON
contract; who produced the JSON --- one vendor's agent, another's, or a human
--- is irrelevant to the ledger.
\end{description}

This paper describes the methodology (\S\ref{sec:method}), its evidence engine
(\S\ref{sec:engine}), the reference implementation (\S\ref{sec:impl}), and the
design-validation evidence accumulated so far (\S\ref{sec:validation}),
followed by limitations (\S\ref{sec:limitations}) and future work
(\S\ref{sec:future}).

\section{Related Work}\label{sec:related}

\textbf{Underspecification.} Ji et al.~\cite{ji2026guessing} show that
underspecified operational instructions cause agents to guess rather than
stop, violating action boundaries in 55.8--67.8\% of runs. spec-loop's
greenfield front is a direct response: the specification is produced through a
structured interrogation in which the agent proposes and the human decides,
and inviolable constraints are captured in a constitution file whose breach is
a first-class blocker, never a trade-off.

\textbf{Building to the test.} Ma et al.~\cite{ma2026building} demonstrate
that agents deliver what is checked, not what is requested, and name the
missing disposition \emph{validation self-awareness}. spec-loop treats this as
its central adversary: acceptance closes per criterion on measured test
evidence; findings cannot be closed without new test material (\emph{find bugs
once}); qualitative verdicts require file-and-line evidence or they do not
count; and the workflow state that authorizes a pull request is derived from
ledger facts, never self-declared by the agent.

\textbf{Process taxonomies.} de Macedo~\cite{demacedo2026prompt} surveys agent
development frameworks along six dimensions (specification, context, roles,
execution, validation, portability) and identifies recurring risks:
specification drift, artifact fragility, and platform dependence. spec-loop
addresses each explicitly: drift through a rebuild test that scores whether
the specification package alone can regenerate the system; fragility through
an integrity-checksummed ledger; and platform dependence through the contract
architecture and a dependency-free engine.

\textbf{Practitioner doctrine.} The methodology's check classification adapts
the computational-versus-inferential distinction discussed in the
practitioner literature on generative AI in software
delivery~\cite{fowler2023exploring}. Architecture decision records follow
Nygard~\cite{nygard2011adr}, extended with machine-checkable implementation
plans. Golden-master capture follows approval-testing
practice~\cite{approvaltests}, with one addition central to the method: the
agent may author the capture harness but is mechanically forbidden --- by a
pre-tool-use hook, not by instruction --- from writing an approved fixture.
The methodology was also systematically cross-examined against classical
doctrine~\cite{thomas2019pragmatic}; \S\ref{sec:validation} reports the ten
resulting improvements. The design stance throughout is an application of
Goodhart's law~\cite{strathern1997improving} to agent verification: any signal
the agent can see, it will eventually optimize.

\section{The Methodology}\label{sec:method}

\subsection{Two specification fronts}

\textbf{Greenfield (discovery).} The human brings an idea; the agent brings
the \emph{shape}. The front proceeds one question at a time, each question
carrying the agent's recommended answer, walking a fixed agenda: purpose,
domain model, operation surface, architecture options with trade-offs, dirty
cases and failure behavior, inviolable constraints, out-of-scope, acceptance
criteria with stable traceable identifiers, a declared quality bar, and
residual risks. High-uncertainty risks trigger a time-boxed \emph{spike} whose
only legitimate output is a decision record with lessons --- spike branches are
mechanically refused at the pull-request gate. The front emits a versioned
specification package: context glossary, constitution, domain model,
specification, decision records with implementation plans and verification
checklists, acceptance file, risks, and handoff.

\textbf{Brownfield (reverse discovery).} For an existing system, the direction
of trust inverts: observable behavior is ground truth, and the agent may
produce \emph{only facts} --- a system map from static analysis and a golden
suite captured mechanically at the boundaries. The agent never authors a
specification of what the old system does or why: a specification written by
the agent about legacy code encodes the same partial reading of the code that
loses behavior silently, which is precisely the blind spot the golden exists
to counter. The human writes the migration specification by reading the facts.

\subsection{The development loop}

The loop consumes the specification package and proceeds through phases, each
guarded: plan (constitution first; no acceptance criteria, no build); a
coverage gate that triggers boundary characterization tests when the safety
net is insufficient; build under decision-record discipline (the agent
consults records before touching governed areas and must propose --- never
silently take --- new architectural decisions); a simplicity gate scoring the
diff against declared budgets, with test code excluded so that writing tests
is never penalized; a severity-gated review loop that \emph{converges instead
of chasing zero} (findings below the declared gate are deferred, not fixed
forever); an integration pass across repositories; late test-writing against
stabilized code; an optional rebuild test measuring specification
completeness; and a pull request opened only when the derived workflow state
says the facts allow it. The loop stops at merge: merging is human.

An escalation contract enumerates the situations in which the agent must stop
and ask --- iteration cap, oscillating findings, a previously-passing test now
failing non-trivially, contradictory tool directives, decisions of
architectural weight, and any constitution breach. Escalations and their
closures are recorded events; an open escalation caps the readiness score
until a human resolves it, and resolving a blocker requires an \emph{escape
analysis}: which gate or test should have caught this, and what was done about
it.

\section{The Evidence Engine}\label{sec:engine}

The engine is a single dependency-free Python file. It never runs an LLM and
never runs the tools it scores: it \emph{ingests} their reports (JUnit-family
XML, Cobertura and lcov coverage, nine linter formats across nine language
stacks) and validates structures. Every verdict is persisted in a
deterministic ledger with an integrity checksum; external mutation or a
truncated write blocks loading with a recovery message.

\textbf{Fact gates (block).} A byte-comparison golden gate (with declared
volatile fields masked \emph{visibly} --- masking is reported, never silent); a
gate-integrity check that blocks diffs that weaken the measuring apparatus
(deleted or disabled tests across all nine stack conventions, lowered
thresholds, added secrets); simplicity budgets over the diff; structural
specification linting (missing sections, zero traceable acceptance criteria,
duplicate identifiers); and a derived state machine: the workflow state
(plan, build, qa, escalated, pr-ready) is \emph{computed} from ledger facts.
A self-declared state machine would be narrated state --- the exact category of
signal the methodology distrusts --- so there is nothing to declare and no
illegal transition to police.

\textbf{Structured guesses (advise).} Prose heuristics over specifications;
regression capture (closing findings without adding a single non-blank test
line is flagged as \emph{narrated closure}); plateau and stop-signal
advisories over the finding history; and a rubric layer: a versioned file of
weighted qualitative criteria with anchor examples, negative criteria, and a
threshold, graded by any runner against a JSON contract with
evidence-or-nothing semantics --- a verdict that affects the score without a
file-and-line citation does not count. All of these advise by default and
gate only when the human declares it in configuration; the declaration
distinction is surfaced everywhere as \emph{requirement (declared) versus
default (kit opinion)}.

\textbf{Readiness.} A weighted 0--100 score reports the state of the result,
never effort. The dominant dimension is measured acceptance --- criteria closed
by green, name-tagged tests --- with checkbox completion demoted to a narrative
dimension. Hard caps override the average (red tests, open blockers,
unresolved escalations), and every biting threshold states its provenance.
The score reports; it does not gate --- the gates are the facts above.

\section{Reference Implementation}\label{sec:impl}

The reference implementation packages the methodology as seven agent skills
(the two fronts, a precision interview for known features, the loop
orchestrator, golden capture, a two-audience system-documentation generator,
and the rubric grader adapter) plus the engine (23 subcommands). Three
properties matter for portability. First, the skills are markdown instruction
files readable by any instruction-following agent; the engine is invoked
identically from any of them. Second, every LLM-judgment interface is a
contract: the rubric grader ships as a neutral prompt usable by any vendor's
agent, a raw API call, or a human filling the JSON by hand --- the
implementation's test suite exercises exactly that path, with no LLM in the
loop. Third, the one mechanical prohibition (the agent must not write approved
golden fixtures) is enforced by a pre-tool-use hook at the runtime boundary,
not by instruction --- during development of the implementation itself the hook
blocked one of its author's own commands, which is the property working as
designed.

The implementation includes an installation-diagnosis subcommand in the
spirit of \texttt{flutter doctor} (interpreter, toolchains per stack, hook
presence and registration, skill integrity, per-project configuration), a
smoke suite of 140 executable checks over a synthetic ledger, and three
installation modes (per-project, per-user, and as a packaged plugin for one
agent runtime --- packaging, not dependency).

\section{Design Validation and Self-Application}\label{sec:validation}

No controlled external evaluation has been conducted yet; we report the
design-validation evidence honestly and completely.

\textbf{Adversarial audits.} Before stabilization, the methodology and its
implementation were subjected to two adversarial audits totaling 231 agent
evaluations: a seven-lens audit (54 agents; findings attacked by skeptics
defaulting to \emph{refuted}) confirmed 33 weaknesses and refuted 13, with the
central finding that the core principle was \emph{inverted in the code} ---
fact gates existed but were not wired to block. A second audit verified 171
claims across documentation, code, and design intent. The response re-wired
every fact gate into the engine and re-passed all documentation against the
implemented reality (130 truthfulness edits), converting the principle from
slogan to enforced property.

\textbf{Doctrine cross-examination.} The methodology was systematically read
against classical software-engineering doctrine~\cite{thomas2019pragmatic},
producing 8 validations, 12 tensions with recorded resolutions, and 10
actionable improvements --- all ten implemented, each as a separate release
with its own regression checks. In the two deepest tensions the methodology
deliberately departed from the letter of the doctrine to defend its own
principle: readiness caps continue to bite by default (their existence is a
definition; only the numbers are opinion, and they now state their
provenance), and the workflow state machine is derived rather than declared.

\textbf{Self-application.} The implementation is developed under its own
process. Across fifteen releases, each engine change carried smoke checks in
the same commit, a fresh-context review before merge, and a five-way version
consistency check enforced by the suite. Fresh-context reviews found real
defects prior to merge in 13 of 14 reviewed releases, including: test-file
classifier gaps that made a documented guarantee an overclaim; a blank-line
loophole in regression-capture evidence; a review-suite check that passed for
the wrong reason (its precondition never held); and silent last-wins
semantics on duplicated grader verdicts --- a grader gaming vector. We take the
consistency of this defect stream as evidence for the methodology's core
premise: independent, fresh-context checking against recorded criteria
catches what the authoring context cannot see.

\section{Limitations and Threats to Validity}\label{sec:limitations}

\textbf{No controlled evaluation.} All evidence is design validation and
self-application by the methodology's author; effect sizes on team
productivity, defect escape rates, or comparison against the frameworks
surveyed in~\cite{demacedo2026prompt} are unmeasured.
\textbf{Single-operator bias.} Self-application shares authorship context;
although reviews run in fresh contexts, they share the model family and the
author's configuration.
\textbf{LLM non-determinism.} The rubric layer's grades vary across runs;
the methodology bounds the blast radius (advisory by default,
evidence-or-nothing, human-declared gating) but does not eliminate variance.
\textbf{Scope.} The loop targets a single operator driving one non-trivial
change; multi-operator concurrency is out of scope, and the ledger is
single-writer by design.
\textbf{Heuristic residue.} Some gates carry documented approximations
(path-based test classification, stale-report windows); each is disclosed at
the point of use, but disclosure is not elimination.

\section{Future Work}\label{sec:future}

Planned next steps are: read-only dry runs of all nine stack adapters against
real repositories; a multi-project dogfooding study with operators other than
the author; a controlled comparison against the frameworks
in~\cite{demacedo2026prompt} along their six taxonomy dimensions, with
particular attention to the validation and portability axes; and an
evaluation of the rubric layer's inter-run variance under anchored versus
unanchored criteria, connecting to the validation self-awareness agenda
of~\cite{ma2026building}.

\section{Conclusion}\label{sec:conclusion}

spec-loop's contribution is not a new agent architecture but a discipline for
placing trust: facts block, guesses advise, the measured overrides the
narrated, and the human decides at exactly three points. The methodology's own
history --- a central principle found inverted by adversarial audit, then
wired into an engine that now blocks its own author when he crosses the line
--- suggests that for LLM-assisted development, the difference between a
methodology and a slogan is whether the checks are executable. The reference
implementation will be released as open source (MIT).

\begin{thebibliography}{9}

\bibitem{ji2026guessing}
Z.~Ji, Z.~Zhang, C.~Xu, Z.~Li, Y.~Gao, S.~Wang, and S.-C.~Cheung,
``Coding Agents Are Guessing: Measuring Action-Boundary Violations in
Underspecified DevOps Instructions,'' arXiv:2607.02294 [cs.SE], 2026.

\bibitem{ma2026building}
Y.~Ma, B.~Kereopa-Yorke, and B.~Schultz,
``Building to the Test: Coding Agents Deliver What You Check, Not What You
Requested,'' arXiv:2606.28430 [cs.SE], 2026.

\bibitem{demacedo2026prompt}
S.~Oliveira de Macedo,
``From Prompt to Process: a Process Taxonomy and Comparative Assessment of
Frameworks Supporting AI Software Development Agents,''
arXiv:2606.04967 [cs.SE], 2026.

\bibitem{thomas2019pragmatic}
D.~Thomas and A.~Hunt,
\emph{The Pragmatic Programmer: Your Journey to Mastery}, 20th Anniversary
Edition. Addison-Wesley, 2019.

\bibitem{nygard2011adr}
M.~Nygard, ``Documenting Architecture Decisions,'' blog post, 2011.
\url{https://cognitect.com/blog/2011/11/15/documenting-architecture-decisions}

\bibitem{fowler2023exploring}
M.~Fowler et al., ``Exploring Generative AI,'' memo series,
martinfowler.com, 2023--2026.
\url{https://martinfowler.com/articles/exploring-gen-ai.html}

\bibitem{approvaltests}
L.~Falco et al., ``ApprovalTests,'' \url{https://approvaltests.com}.

\bibitem{strathern1997improving}
M.~Strathern, ``\,`Improving ratings': audit in the British University
system,'' \emph{European Review}, vol.~5, no.~3, pp.~305--321, 1997.

\end{thebibliography}

\end{document}
